\documentclass{article}

%\usepackage{loubookman}
\usepackage{loustandard}
\usepackage{graphics}
\usepackage[tt]{codedocn}


\newcommand{\OsiTwo}{OSI2\xspace}
\newcommand{\OsiTwoTop}{\pgmid{Osi2::API}\xspace}
\newcommand{\OsiTwoControl}{\pgmid{Osi2::ControlAPI}\xspace}
\newcommand{\OsiTwoOsi}{\pgmid{Osi2::Osi1API}\xspace}
\newcommand{\OsiTwoModel}{\pgmid{Osi2::ModelAPI}\xspace}
\newcommand{\OsiTwoSolve}{\pgmid{Osi2::SolveAPI}\xspace}
\newcommand{\revisionDate}{November 10, 2011}

\begin{document}
\markright{\OsiTwo \hfil Design and Architecture \hfil \revisionDate}

\heading{\OsiTwo Design and Architecture}

This document describes the design and architecture plans for \OsiTwo.

\OsiTwo defines application programming interfaces (APIs) for common
optimisation tasks (maintaining a model, obtaining an optimal solution to
the model, \etc).
Plugin libraries supply objects to implement specific APIs.
Libraries can be dynamically loaded and unloaded by a plugin manager.
When a library is loaded, it registers the set of APIs that it supports
with the plugin manager.

Clients use the \OsiTwoControl interface to load libraries and create
objects that implement the APIs supported by the libraries.
The same interface is used to destroy objects and unload libraries.

In response to a request to load a plugin library, the \OsiTwoControl object
interacts with the plugin manager to load and initialise the library.
In response to a request for an object implementing a specific API or
combination of APIs,
the \OsiTwoControl object interacts with the plugin manager to invoke the
appropriate factory method to produce an object supporting the requested
API(s) and hands this
object over to the client.
Clients interact with these objects according to the interface defined
by the API\@.
The \OsiTwoTop class serves as the base class for all APIs.

There are a number of overarching design goals:
\begin{itemize}
  \item
  Dynamic loading of solvers and other functional modules to implement
  specific APIs.

  \item
  A lightweight and transparent plugin manager implementation.

  \item
  Small, focused APIs, so that
  individual APIs are kept to a
  manageable size and implementation of an object in support of
  an API is a manageable task.

  \item
  Support for interaction between API objects to provide complex
  functionality: at minimum, composition and replacement.
\end{itemize}

\subheading{API Architecture}

\OsiTwo APIs derive from a common top-level virtual class \OsiTwoTop as
shown in Figure~\ref{fig:OsiTwoAPIHierarchy}.
\begin{figure}[b]
\begin{center}
\includegraphics{apihierarchy}
\end{center}
\caption{Class Hierarchy for Osi2 APIs}
\label{fig:OsiTwoAPIHierarchy}
\end{figure}

At present, \OsiTwoTop provides support for a single capability: object
identification, in the form of storage and retrieval of a pointer to an
opaque data block.

\subheading[l]{\OsiTwoControl}

The \OsiTwoControl API provides a simple user interface for loading and
unloading plugin libraries and for creating and destroying objects.

The methods that support plugin library loading (\pgmid{load}) associate a
`short name' string with a plugin library.
Subsequent references to the library for object creation or unloading use
only the short name.
A default name for the plugin library can be constructed from the short name
or the full name can be specified by the user.
The full path to the plugin library is constructed using a default path or a
path specified by the user.
All interaction with the plugin manager is hidden from the client.

To create an object (\pgmid{createObject}),
the client specifies the API(s) that the object should
support using the name(s) registered with the plugin manager.
An optional parameter allows the user to specify that the object must be
created by a particular plugin library.

To destroy an object (\pgmid{destroyObject}),
the client specifies the object to be destroyed.
It is possible to use the standard C++ \pgmid{delete} operator to destroy an
object, but clients should use the \pgmid{destroyObject} method
to ensure that any additional actions required by the plugin library are
performed.

To unload a library (\pgmid{unload}), the client specifies the short name
of the library.

Because one library's implementation of an API can differ from another
library's implementation of the same API, it is necessary that destruction of
an object be performed by the same library that created it.
The \OsiTwoControl implementation uses the identification hook in \OsiTwoTop to
store identification information when the object is created.
The client need only provide a pointer to the object when asking for it to be
destroyed.

\subheading[l]{\OsiTwoOsi}

As a proof-of-concept for the \OsiTwo architecture, an \OsiTwoOsi API has been
implemented.
This API allows any implementation of the \pgmid{OsiSolverInterface} API to
be used as a plugin in the \OsiTwo architecture.
\OsiTwoOsi1 is declared as
\begin{verbatim}
class Osi1API : public API { ... }
\end{verbatim}
and declares all \pgmid{OsiSolverInterface} methods to be pure virtual.

Shims have been implemented for the clp and glpk solvers using a multiple
inheritance strategy.
Using the clp shim as an example, the declaration is
\begin{verbatim}
class Osi1API_ClpHeavy : public Osi1API, public OsiClpSolverInterface { ... }
\end{verbatim}
A typical implementation of a method looks like
\begin{verbatim}
inline void initialSolve() { OsiClpSolverInterface::initialSolve() ; }
\end{verbatim}
The only feature requiring nontrivial adaptation is the
\pgmid{ApplyCutsReturnCode} structure returned by \pgmid{applyCuts}.
This is handled with an adapter method as:
\begin{verbatim}
inline Osi1API::ApplyCutsReturnCode applyCutsPrivate(const OsiCuts &cuts,
					      double eff = 0.0)
{ 
  const OsiSolverInterface::ApplyCutsReturnCode &tmp =
      OsiClpSolverInterface::applyCuts(cuts,eff) ;
  const Osi1API::ApplyCutsReturnCode
      retval(tmp.getNumInconsistent(),
	     tmp.getNumInconsistentWrtIntegerModel(),
	     tmp.getNumInfeasible(),
	     tmp.getNumIneffective(),
	     tmp.getNumApplied()) ;
  return (retval) ;
}
\end{verbatim}
The declaration of \pgmid{ApplyCutsReturnCode} is duplicated in
\pgmid{Osi1API} with the addition of a constructor that allows all data values
to be specified at construction.
Then \pgmid{applyCuts} is implemented as
\begin{verbatim}
virtual ApplyCutsReturnCode applyCuts(const OsiCuts & cs,
                                      double effectivenessLb = 0.0)
{ 
  return (applyCutsPrivate(cs,effectivenessLb)) ;
}
\end{verbatim}
The client includes \pgmid{Osi2Osi1API.hpp} and works entirely with the
\OsiTwoOsi API.
In all other respects, it's exactly as if the client was working with the
original \pgmid{OsiSolverInterface} API.

\subheading[l]{The Frontier}

Further work on \OsiTwo requires some fundamental decisions about the type of
usage that should be supported.
As a running example,
assume that \OsiTwo defines an \OsiTwoModel API for model construction and
modification and an \OsiTwoSolve API for `solving' the model by finding an
optimal solution.
Consider the following scenario:
The client starts out constructing a model using an object which supports the
\OsiTwoModel API\@.
Once the model is constructed, the client would like an object which supports
the \OsiTwoSolve API to be used to solve the model just constructed.
There may be multiple iterations of modify and solve, and the user will expect
that changes to the model are seen by the solver.

Given an object implementing some set of APIs, the client might want to add to
that set or reduce it.
\begin{itemize}
  \item
  A reduction would specify that support for a particular set of APIs is no
  longer required.
  This might be necessary in order to add an API that is incompatible with the
  API being removed.

  \item
  An addition would specify that support for a particular set of APIs be added
  to the object.
\end{itemize}

The client needs to be able to choose whether the object resulting from an
API change is a new, independent object, or the same object with a different
set of capabilities.
In \OsiTwo, the design decision is that the default should be to alter
the existing object whenever possible, returning a new object only when
absolutely necessary.

In the context of the running example, should the object that supports the
\OsiTwoModel API be independent from the object that also supports the
\OsiTwoSolve API\@?
If the model is initially constructed using the model construction
capabilities of a solver, the existing object will do just fine.
If the model is initially constructed using a standalone model representation
optimised for efficient modification, it may be necessary to load the
model into a solver object in order to provide the solve capability.

\OsiTwo should provide the ability to explicitly request the creation of an
independent object.
It should also provide the capability to inquire whether the desired change in
APIs will force the creation of an independent object.

Should the client be expected to work with two distinct (but coordinated)
objects, one supporting the \OsiTwoModel API and one supporting the
\OsiTwoSolve API\@?
If not, whose responsibility is it to define a single class that supports
both APIs?

In \OsiTwo, if a single object supporting multiple APIs is what is desired,
it will be the client's responsibility to derive the necessary class
from \OsiTwoTop.
The client knows the exact set of combinations that will be needed in the
application.
\OsiTwo could guess at the set of desirable combinations and declare
corresponding classes, but combinatorial
explosion will defeat any attempt to exhaustively enumerate the possible
combinations.

How will the plugin library support interactions between API objects?
Independent of the client's view, should the implementation of this composite
capability be two interacting objects or a single object providing both
APIs?
Here are two possible models:
\begin{itemize}
  \item
  Multiple interacting plugin objects can be supported within the \OsiTwo APIs.
  One can divide the objects into two broad classes: data objects and analytic
  (action) objects.
  In the running example, we have a single data object (specified by the
  \OsiTwoModel API) and a
  single analytic object (specified by the \OsiTwoSolve API).
  The general case is a many-to-many interaction.
  An analytic object can process multiple data objects.
  A data object can be processed by multiple analytic objects.

  As a general algorithm for synchronisation, data objects could keep sequence
  numbers that track modifications to the data object.
  Each analytic object could record the current sequence number when it
  updates itself from its associated data objects.
  At the point where an analytic object is asked to process its associated
  data objects, it queries each data object to determine its current sequence
  number and updates itself as necessary.
  Updates are of most concern where the analytic object maintains an internal
  copy of the data object.
  (True, for instance, for most solvers, which load in a problem and maintain
  an internal representation that is suited to the needs of the solver's
  algorithms.)

  Data objects could also maintain a list of modifications.
  This would allow for incremental update of analytic objects, where that was
  more efficient than bulk update.
  It also provides a mechanism for an undo/redo capability.

  To interface to the single-object client-facing model, the client-facing class
  would simply maintain pointers to the underlying objects.

  \item
  Single (monolithic) plugin objects implementing some fixed combination of
  APIs can be supported entirely within the plugin library.
  Coordination is trivial in the sense that it is an internal responsibility
  of the object.
  No additional support is required within the \OsiTwo API hierarchy.

  As with the multiple interacting object model, the client-facing class
  need only maintain a (single) pointer to the appropriate object implemented
  by the plugin library.

  A limitation of this model, as opposed to the multiple interacting object
  model, is that arbitrary composition is not possible.
  The plugin developer faces the same combinatorial explosion
  mentioned earlier.

  An advantage of this model is that it captures a large class of analytic
  tools (most solvers, for example) which provide extensive APIs for
  building, manipulating, and solving models.
  Referring to the running example, nearly all solvers will implement both the
  \OsiTwoModel API and the \OsiTwoSolve API (a consequence of maintaining an
  internal representation tailored to the algorithms implemented in the solver).
  `Augmenting' the \OsiTwoModel API with the \OsiTwoSolve API is a
  trivial operation, because the object provided by the plugin library always
  supports both APIs.
\end{itemize}

One can ask whether there is a convincing use case that justifies the
complexity of the multiple interacting objects model.
Put another way, why do the majority of existing tools fit the single object
model\@?
A plausible explanation is that a tool such as a solver is expected to be 
complete in itself.
Given the infrastructure, would developers target the multiple object model,
where the ability to use plugin modules would reduce the need to reimplement
common functions and thus lower the initial effort to implement an
algorithm\@?
Or would they be deterred by the loss of control and possible loss of
efficiency implied by dependence on an external plugin object?

Another issue is the nature of the interface between the \OsiTwo object
and the implementing object provided by the plugin library.
The client-facing interface of \OsiTwo objects is defined by the \OsiTwo APIs,
but the plugin-facing interface need not be the same.
Arguably, the plugin-facing interface should be parsimonious, particularly in
the monolithic plugin object model, where each additional capability is
added on top of previous capabilities.
The danger is that the work required to implement a minimally functional
plugin will become an insurmountable bar (one of the criticisms labelled at
the existing \pgmid{OsiSolverInterface}).
This must be balanced against efficiency: if the set of methods is too small
and/or primitive, too many basic operations will be required to implement a
complex operation.


\subheading{Plugin Manager Architecture}

The plugin manager architecture used in \OsiTwo is broadly adopted from an
architecture originally described by G.~Sayfan \cite{say07}.
As plugin libraries are loaded, the plugin manager registers the capabilities
(APIs) provided by the libraries.
When a client asks for an object to implement an API, the plugin manager
locates a library which can satisfy the request and invokes a factory method
to create the object.
APIs are identified by strings.
Maps are used to associate an API identification string with the information
necessary to create an object that implements that API.

There is one (static) instance of the plugin manager, returned by a call to
\pgmid{getInstance}.
The plugin manager provides services and information to plugins.
Currently, the plugin manager provides a registration service for APIs.
The plugin manager also implements the capability to provide general-purpose
services\footnote{%
  No general-purpose services are implemented at present.}.

A deficiency in the original architecture is that it assumes a very structured
interaction between plugin libraries and the client program.
\begin{itemize}
  \item
  There is no provision for maintaining state to manage a library or individual
  APIs supported by the library.

  The plugin libraries used for \OsiTwo are in general not stateless.
  The interface between the plugin manager and the library  has been expanded
  to allow a reference to a management object to be returned to the plugin
  manager during library initialisation and passed back to the plugin library
  whenever the library is asked to create an object to implement an API.

  \item
  There is no provision to associate registered APIs with a
  particular plugin library, hence there is no way to unload a single plugin
  library.

  The \OsiTwo architecture supports loading and unloading of individual plugin
  libraries throughout execution.
  It also allows a client to specify that a particular plugin library should
  be used to satisfy a request to create an object implementing some API\@.
  To support these capabilities, the plugin manager maintains
  bookkeeping information to associate registered APIs with a specific library.
\end{itemize}

The working sequence to load a plugin library, create an object, destroy an
object, and unload the library, is described below.
The first two steps (loading the shared library, initialising the plugin and
registering APIs) are encapsulated in \pgmid{PluginManager::loadOneLib}, so
that load and initialisation of a plugin library is an atomic action from the
point of view of a client.
\begin{enumerate}
  \item
  Load the plugin library.

  This is accomplished by a call to the static \pgmid{load}
  method in the \pgmid{DynamicLibrary} class, which will invoke the system
  \pgmid{dlopen} method\footnote{%
      Support for Windows will be provided by has not been tested.}.
  The plugin manager supports specification of a default plugin
  directory\footnote{%
    We need to add infrastructure to check the environment. We might also
    want to allow for a search list of directories.
    A design decision is whether this capability should move to
    \OsiTwoControl or be replicated in \OsiTwoControl.}.

  \item
  Initialise the plugin library and register APIs implemented by the library.

  Given a \pgmid{DynamicLibrary} object, \pgmid{getSymbol} is used to
  retrieve the plugin library's initialisation (\pgmid{InitFunc}) method,
  which must be named \pgmid{initPlugin}.
  This method is then invoked,
  passing a \pgmid{PlatformServices} structure as a parameter.
  The \pgmid{InitFunc} is responsible
  for getting the plugin library ready to work and registering
  its capabilities (APIs).

  For each API supported by the plugin library, a \pgmid{RegisterParams} object
  is created specifying the identifying string for the API, create
  (\pgmid{CreateFunc}) and destroy (\pgmid{DestroyFunc}) methods,
  descriptive information (version and language), and an associated
  object to hold any state information the library wishes to maintain for
  objects supporting this API.
  The \pgmid{RegisterParams} object is handed back to the plugin manager
  as a parameter to the registration (\pgmid{RegisterFunc}) method
  supplied in the \pgmid{PlatformServices} object.
  The \pgmid{RegisterFunc} method adds the \pgmid{RegisterParams} object to
  an appropriate temporary structure in the plugin manager
  according to the API identity string (exact match or wildcard).

  Once all APIs are registered, the \pgmid{InitFunc} method stores a pointer
  to the library's control object in the \pgmid{PlatformServices} structure
  and returns to the plugin manager.
  The return value is a pointer to the library's cleanup (\pgmid{ExitFunc})
  method.

  At this point it is no longer possible for initialisation of the plugin
  library to fail.
  The plugin manager copies the API registration information from the
  temporary structures to permanent structures to complete the load of the
  library.

  \item
  Ask the plugin manager for an object that implements an API.

  Using the plugin manager's \pgmid{createObject} method, the client asks for
  an object that implements a particular API by specifying the identity
  string registered for that API\@.
  The \pgmid{createObject} method searches the maps held by the plugin manager
  looking for a suitable match.
  If a match is found, 
  an \pgmid{ObjectParams} object is created, including the API requested, a
  \pgmid{PlatformServices} object, and a pointer to the control object for
  the API.
  (A pointer to the control object for the library is passed in the
  \pgmid{PlatformServices} object.)
  The \pgmid{ObjectParams} object is passed as a parameter to the
  \pgmid{CreateFunc} that the plugin registered for this type of API.
  The \pgmid{CreateFunc} returns a pointer to an object that implements the
  requested API.

  The client can specify that the API must be produced by a specific plugin
  library.
  This allows full client control when more than one library can produce an
  object that implements an API.

  \item
  Ask the plugin manager to destroy an object.

  It is possible to simply invoke \pgmid{delete} on an object, but use of the
  plugin manager's \pgmid{destroyObject} method allows some additional
  functionality.
  The \pgmid{destroyObject} method invokes the plugin library's
  \pgmid{DestroyFunc} for the object, passing the same \pgmid{ObjectParams}
  parameter block used for \pgmid{createObject}.
  This allows the plugin library to perform additional management at the API
  and library levels, if necessary.

  The plugin manager requires that the request to destroy the object specify
  the plugin library that produced the object.
  This is necessary to ensure that the
  request to destroy an object goes to the plugin library that created
  it\footnote{%
  The plugin manager does \textit{not} track objects; it's purely an
  intermediary.
  Given that objects can be deleted using the standard \pgmid{delete}
  operator, there's no point in even trying.
  While the \pgmid{destroyObject} method requires that the client specify a
  plugin library, the plugin manager cannot ensure that the library specified
  is in fact the library that produced the object.
  It's the client's responsibility to keep this straight.}.

  \item
  Ask the plugin manager to unload a library.

  The \pgmid{unloadOneLib} method will scan the plugin manager's API
  registration maps and remove all entries associated with the library to be
  removed.
  It will then invoke the library's \pgmid{ExitFunc}.
  Finally, the destructor for the dynamic library object will invoke
  \pgmid{dlclose} to unload the library.
\end{enumerate}

A given plugin library can supply multiple APIs, and a client can ask for an
object that implements some combination of those APIs by specifying multiple
API identity strings in a request to create an object.
Any given request to a plugin library can return only one object.
It's the business of the plugin library to keep track of what particular APIs
or combination of APIs it supports, and what objects it's given out, to the
extent that this is necessary.

There is no distinct `capability inquiry' method for the plugin
manager\footnote{%
    At the time of writing, it's looking like this would be a useful
    extension.}.
Given a string specifying one or more APIs to be supported by an object, the
plugin manager first looks to see if some plugin library has registered
exactly that API string.
If not, the plugin manager looks to see if any libraries have registered
wildcard capability.
For each such plugin library, the plugin manager sends the API string to the
library.
If the library returns an object, the plugin manager registers the API string
so that future requests will be satisfied with an exact match.

\subheading[l]{Shims}

The \OsiTwo term for an implementation of a plugin library is a `shim'.
Two shim architectures have been tried, called `heavy' and `light'.

A `heavy' shim includes any underlying libraries in the link that creates the
shared library for the shim.
A `light' shim dynamically loads additional libraries as part of its
initialisation.

As an example, there are light and heavy shims for the clp solver.
The heavy shim links includes \pgmid{libClp} as part of the link.
The light shim dynamically loads \pgmid{libClp} as part of its
initialisation\footnote{%
    At present, this is accomplished with direct calls to methods in the
    \pgmid{DynamicLibrary} class.
    It's not yet clear whether this is the appropriate choice.
    The alternative is to use the plugin manager, but one of the major
    services provided by
    the plugin manager (registration of APIs) may not be needed here.
    If API registration is used, the relevant code in the plugin manager
    will need to be made reentrant.
    Currently, initialising a library from within the initialisation of
    another library will pollute the temporary structures used to hold API
    registration.}.
The primary advantage to the heavy shim is that C++ classes can be used as
normal.
The disadvantage is that the shim cannot be distributed independently of the
underlying support libraries.

The primary advantage to the light shim is that the shim can be distributed
independently of the underlying libraries.
The disadvantage is that all library methods must be explicitly loaded with
\pgmid{getSymbol} (\pgmid{dlsym}) and are subject to all the restrictions of
`C' linkage.

The \OsiTwoOsi implementations described earlier use a multiple inheritance
architecture, as
\begin{verbatim}
class Osi1API_ClpHeavy : public Osi1API, public OsiClpSolverInterface { ... }
\end{verbatim}
with methods implemented as
\begin{verbatim}
inline void initialSolve() { OsiClpSolverInterface::initialSolve() ; }
\end{verbatim}
One can also work with a pointer to the underlying solver object, as
\begin{verbatim}
#include "OsiClpSolverInterface.hpp"

class Osi1API_ClpHeavy : public Osi1API {

  OsiClpSolverInterface *clp_ ;

}
\end{verbatim}
with methods implemented as
\begin{verbatim}
inline void initialSolve() { clp_->initialSolve() ; }
\end{verbatim}
At the expense of an additional indirection, this approach provides the
ability to change the solver object.
In a more general context, this approach offers considerable flexibility to
assemble the methods required for an API from multiple underlying objects with
fewer restrictions than multiple inheritance.

\subheading[l]{The Frontier}

What is the correct model for providing services to a plugin library?
Here are three models that appear to be useful:
\begin{itemize}
  \item
  Use the existing service function hook, \pgmid{InvokeServiceFunc}, provided
  by the plugin manager.
  This seems the appropriate route for stateless services, particularly
  services where efficiency isn't a dominant concern.
  Logging would be an example of such a service.

  \item
  A plugin library could request from the plugin manager an object that
  implements a particular API\@.
  These might be supplied by some other plugin library, or they could be
  `static' plugins compiled into the plugin manager.
  This would be a good model for services that require state, or where
  there will be frequent use and efficiency is a concern.
  Better to hand out an object that can be used repeatedly without further
  interaction with the plugin manager.

  \item
  A service can be implemented in an API\@.
  For some services, it will be necessary to insert functionality into
  calls from the client to the plugin.
  This can be handled by placing implementation code in the API definition
  class, where it can execute before or after invoking methods provided by the
  plugin library.

  Caching can serve as an example.
  A caching service that maintains a data structure embedded in the API
  could minimise the number of times
  that it is necessary to call the underlying plugin library.
  This requires code in the API to manage the embedded data structure.
  A ModelAPI, for example, could maintain an embedded representation
  of the model which it could modify incrementally in response to calls from
  the client, completely avoiding calls to the underlying library.
  At some point, the client would request that the underlying library process
  the model.
  The caching code could push the revised model to the underlying library
  prior to invoking the library's processing method.
  (To couch this in terms of specific calls, consider a solver plugin.
  An API could maintain a cached copy of a constraint system, trapping
  modifications within the API.
  When the client invokes an optimisation method, the API would note that
  the model was modified and push the updated model to the solver before
  invoking the solver's optimisation method.)
\end{itemize}


\bibliographystyle{plain}
\bibliography{osi2design}

\end{document}

